\section{標準六項への補完を次に使うための要点}

今回の主対象は $\Lambda=R_{(0,1,2,3,4,5)}$ と $B=e\Lambda e$ の組である。
可換ファイバーの標準埋め込みを固定する。
次の三つを別々の問題として扱う。
\begin{enumerate}
\item 一次方向を周囲へ持ち上げること。これは $\rho^2$ の同型により常に可能。
\item 一次方向を固定し、補完可能な全形式族を何か一つ作ること。必要十分条件はPoisson性。
\item 高次の乗法まで指定済みの全形式 $B$ 変形を補完すること。一次Poisson性だけでは足りない。
\end{enumerate}
第二区分の十分性は発散零代表と量子化による。必要性は二次のHochschild障害による。
第三区分については、下の二変数族が全Artin基底での正確な試金石になる。

\subsection{二次的基底の認定方法}
六項のminiversal基底と四項の基底は、compatibleな選択のもとで
\[
 R_\Lambda\simeq k[[u_1,\ldots,u_{53}]]/(q_1,\ldots,q_{38}),\qquad
 R_B\simeq k[[u_1,\ldots,u_{53}]]
\]
となり、制限は基底上の閉埋め込みになる。
三十八式はJacobi式の独立な同次二次係数である。
基底の座標は標準的でなく、hullは自己同型を持つ変形関手のfine表現ではない。

この結果の認定補題は他の問題にも再利用できる。
同次二次生成イデアル $J\subset P=k[[u]]$ の零点全体上に接写像が同型の全形式族があれば、
miniversal基底 $P/I$ からの写像を取り、形式座標変更により $I\subset J$ とできる。
さらに二次障害で $I_2=J_2$ が分かれば
$J\cap\mathfrak m^3=\mathfrak mJ$ とNakayamaから $I=J$ となる。
高次の逐次消去を続ける代わりに、二次軌跡全体の全形式実現を探すことが有効である。

\subsection{全Artin基底での二変数モデル}
$[x_1,x_2]=2s x_1x_3$、$x_3$ 中心のOre部分と、
$x_iw_1=(1+t\delta_{i2})v_i$、$w_0x_i=(1-t\delta_{i2})v_i$ という平方零部分から
$B_{s,t}$ を構成する。補完存在条件は正確に $st=0$ である。
これは $s,t$ が独立な形式変数の場合だけでなく、任意の局所Artin環における値にも成立する。

必要性の証明で再利用すべき道具は二つある。
第一に、$HH^1$ 制限の全射性と $HH^2$ 制限の単射性から、corner同型を周囲の同型へ持ち上げられる。
第二に、任意の $st\neq0$ に対して
\[
 J=T(st),\qquad T'=T/\mathfrak m_TJ,
 \qquad I=J/\mathfrak m_TJ\neq0
\]
を取ると、$\mathfrak m_{T'}I=0$ の小拡大に還元できる。
そこで純粋な $s$ 項・$t$ 項は軸上の結合律で消え、三次以上の混合項も消えるため、
唯一残る障害は $(st)[\alpha_s,\alpha_t]$ である。
Ore表示の失敗だけでなく、全ての周囲の変形に対する不可能性を証明するにはこの還元が必要である。

$B_{\hbar,\hbar^r}$ ($r\ge2$) は非零Poisson一次方向を持ち、
$\hbar^{r+1}$ を法として補完可能、$\hbar^{r+2}$ を法として補完不能である。
従って、特定次数まで補完できたという計算結果だけから全次数の補完を推測しない。
ただしこれは有限本の方程式の不存在ではない。ここでは全条件が一式 $st=0$ である。

\subsection{今回の原稿構成の判断}
公表稿VIIIは、三次制限・二次的hull・非純次元性・節点モデルを一つの補完問題としてまとめる。
公表稿Iの標準部分代数の一意性とLO比較とは、主定理も必要な背景も異なるため独立稿とする。
既存資料にある四項側の無障害性は、VIIIで短い自己完結した証明として再構成した。
I--VIIの証明範囲は今回変更していない。
量子中間点の随伴次数同定を元のLO代数の実現へ移さないこと、
標準六項の補完不能性を全ての導来埋め込みの不存在へ移さないことを維持する。
